\cotefile{43}
\scanpage{178}
% archive sheet number: 177

\diagram{\begin{tikzcd}[ampersand replacement=\&,column sep=small,row sep=small]
\& \& \& \mathbb Z/p\mathbb Z \arrow[dr,dash,"\mathbb Z/\ell\mathbb Z"] \& \\
\& \& (\unsure{G_a|\mathbb G_a},\alpha_p) \arrow[ur,dash] \arrow[dr,dash,"(\mathbb Z/\ell\mathbb Z)"] \& \& (\mathbb Z/p\mathbb Z) \\
\& (\mu_p) \arrow[ur,dash] \arrow[dr,dash,"(\mathbb Z/\ell\mathbb Z)"] \& \deleted{\unsure{\mathbb V|G_a}} \& (G_a,\alpha_p) \arrow[ur,dash] \& \deleted{\unsure{\mathbb Z/p\mathbb Z}} \\
\uline{(\unsure{G_m|\mathbb G_m})} \arrow[ur,dash] \arrow[dr,dash,"(\mathbb Z/\ell\mathbb Z)"'] \& \& (\mu_p) \arrow[ur,dash] \& \& \\
\& \deleted{\unsure{(G_a,\alpha_p)}} \arrow[ur,dash] \& \& \&
\arrow[from=4-1,to=2-3,bend left=55,dash,"{\text{\unsure{k.connexes|b.connexes}}}"]
\arrow[from=4-1,to=3-2,bend left=60,dash,"{\text{\unsure{multiplicatif connexe}}}"]
\arrow[from=1-4,to=2-5,bend left=70,dash,"{\text{\unsure{fini séparable}}}"]
\arrow[from=4-1,to=5-2,bend right=40,dash,"{\text{\unsure{multiplicatif séparable}}}"']
\arrow[from=4-1,to=4-3,bend right=75,dash,"{\text{\uuline{multiplicatif}}}"']
\arrow[from=4-4,to=3-5,bend right=60,dash,"{\text{\unsure{unipotent k.connexe}}}"']
\arrow[from=4-4,to=2-5,bend right=80,dash,"{\text{\uuline{unipotent}}}"']
\end{tikzcd}}

% Vertical separator between the two upper sketches.
\figure{\begin{tikzpicture}\draw (0,0)--(0,3.2);\end{tikzpicture}}

\diagram{\begin{tikzcd}[ampersand replacement=\&,column sep=small,row sep=small]
\& \& \mathbb Z/p\mathbb Z \arrow[dr,dash,"(\deleted{\mathbb Z/\ell\mathbb Z})"] \& \\
\& (\alpha_p,\unsure{\mathbb V|V}) \arrow[ur,dash] \arrow[dr,dash,"(\deleted{\mathbb Z/\ell\mathbb Z})"] \& \& \uline{(\mathbb Z)} \\
(\mu_p) \arrow[ur,dash] \arrow[dr,dash,"\mathbb Z/\ell\mathbb Z"'] \& \deleted{\mathbb Z/\ell\mathbb Z} \& (\deleted{\mu_p},\mathbb V) \arrow[ur,dash,"(\mathbb Z/p\mathbb Z)"'] \& \\
\& (\mu_p) \arrow[ur,dash] \& \&
\arrow[from=3-1,to=2-2,bend left=65,dash,"{\text{\uline{\unsure{infinitésimal}}}}"]
\arrow[from=2-2,to=2-4,bend left=70,dash,"{\text{\uuline{\unsure{discret sép.}}}}"]
\arrow[from=3-1,to=3-3,bend right=70,dash,"{\text{\deleted{\unsure{infinitésimal}} \unsure{sans torsion additive}}}"']
\arrow[from=3-3,to=2-4,bend right=50,dash,"{\text{\deleted{discret} \deleted{(\unsure{séparable})}}}"']
\arrow[from=3-3,to=2-4,bend right=75,dash,"{\text{\unsure{sans torsion multiplicative}}}"']
\end{tikzcd}}

Supposons la partie \deleted{\unsure{réduite}} \unsure{séparable unipotente connexe déf/k.} \unsure{Alors le diagramme se complète} :

\diagram{\begin{tikzcd}[ampersand replacement=\&,column sep=small,row sep=small]
\& \& \& (\mathbb Z/p\mathbb Z) \arrow[dr,dash,"(\mathbb Z/\ell\mathbb Z)"] \& \\
\& \& (G_a) \arrow[ur,dash,"(\alpha_p)"] \arrow[dr,dash,"(\mathbb Z/\ell\mathbb Z)"] \& \& \uline{(\mathbb Z/p\mathbb Z)} \\
\& (\mu_p) \arrow[ur,dash] \arrow[dr,dash,"(\deleted{\mathbb Z/\ell\mathbb Z})"] \& \& (G_a) \arrow[ur,dash,"(\alpha_p)"'] \& \\
\uline{(G_m)} \arrow[ur,dash] \arrow[dr,dash,"(\mathbb Z/\ell\mathbb Z)"'] \& \& (\mu_p) \arrow[ur,dash] \& \&
\arrow[from=2-3,to=2-5,bend left=65,dash,"{\text{\unsure{fini}}}"]
\end{tikzcd}}

\margin{Si de plus la partie \unsure{séparable de $G_0$ est définie} sur $k$, \deleted{\unsure{nous avons}} \inserted{\unsure{(ce}} \unsure{revient au $=$ de dire} \unsure{que celle de $G_0/\{G_m\}$} \unsure{l'est) alors on a un diagramme} \deleted{\unsure{moins}} plus \unsure{complet}.}

\[
\underbrace{\underbrace{\mu_p,(G_a,\alpha_p),\mathbb Z/p\mathbb Z}_{\vphantom{a}}\,,\,
\underbrace{\mu_p,(\alpha_p,\mathbb V),\mathbb Z/p\mathbb Z}_{\vphantom{a}}}_{\text{\unsure{extension triviale}}}
\]

\deleted{\unsure{$\mu$}}
\diagram{\begin{tikzcd}[ampersand replacement=\&,column sep=large,row sep=large]
\deleted{\unsure{G_m}} \&
\left[(\overbrace{\mu_p),(G_a,\alpha_p}^{\vphantom{A}},\mathbb V),(\mathbb Z/p\mathbb Z)\right]
\arrow[dl,dash,"(\mathbb Z/\ell\mathbb Z)"']
\arrow[dr,dash,"\mathbb Z/\ell\mathbb Z"]
\& (\mathbb Z) \arrow[dl,dash] \\
(G_m) \arrow[rr,dash] \& [\quad] \& {}
\end{tikzcd}}
\[
\underbrace{(G_a,\alpha_p,\mathbb V),(\mathbb Z/p\mathbb Z)}_{\text{\unsure{Ext. un. triviale}}}
\]

\uline{Multiplicatif \unsure{mult.connexe|ind.connexe}} (\unsure{pro-tore})
\[
\left(\begin{array}{c}\text{\unsure{pro-étale}}\\ \text{\unsure{premier à}}\ p\end{array}
\times \text{\unsure{gr. pro-unipotents}}\right)
\]
\margin{\uline{\unsure{pro-discret}} \uline{\unsure{séparable}} \uline{sans torsion}}
\scanpage{180}
% archive sheet number: 179

\[
\mu_p+(G_a,\alpha_p)
\qquad\qquad \mu_p\quad\alpha_p\quad\unsure{\mathbb Z/p|\mathbb Z_p}
\]
\[
(\alpha_p,\unsure{\mathbb V|V})+\mathbb Z/p\mathbb Z
\]
\diagram{\begin{tikzcd}[ampersand replacement=\&]
G_a \arrow[rr,leftrightarrow,"\alpha_p"'] \& \& V
\end{tikzcd}}
\[ ( \]
\[\mu_p\times\text{\unsure{Unipt}},\quad\mathbb V\]

\unsure{Multiplicative}\qquad unipotents\qquad discret \unsure{ordinaire}

\diagram{\begin{tikzcd}[ampersand replacement=\&,column sep=large,row sep=large]
\& G_0 \& \& \\
\text{\unsure{multiplicative} (connexe séparable} \arrow[ur,dash,"{\text{(\unsure{rédt})}}"]
\arrow[dr,dash,"{G_{0\unsure{\mathrm{réd}|\mathrm{rid}}}}"']
\arrow[rru,dash,"{\text{multiplicatif fini radiciel}}"]
\arrow[rrd,dash,"{\text{mult. fini sép.}}"']
\& \& \& {} \\
\& G_{\mathrm{réd}} \& {}
\arrow[ru,dash,"{\text{multiplicatif radiciel fini}}"'] \&
\arrow[from=1-3,to=2-4,dash,"{\text{multiplicatif fini sép.}}"]
\end{tikzcd}}
\deleted{\illegible{1}}

$G_{\mathrm{réd}}$ est un \unsure{sous-gr} si $G$ est multiplicatif, mais pas s'il est unipotent.

$G_0$ est défini sur $K$ \deleted{si $G$ est multiplicatif} \inserted{en tous cas}
\deleted{mais pas} [\unsure{faut qu'une variété} \unsure{alg. connexe} \inserted{\unsure{avec|définie}} \unsure{d'un|dim.} \deleted{\unsure{ait}} \unsure{pt rat /k} ???]

\unsure{unipotent connexe} (\unsure{gr. unip. sép.|gr. irr. sép.})\qquad \unsure{fini} \uline{unipotent} \uline{séparable}

\deleted{\unsure{$G$}} \deleted{\unsure{$M$}}
\[(\unsure{G_m|\mathbb G_m})\quad(\mu_p)(\mathbb Z/p\mathbb Z)\quad\text{Unip. conn.}\]

\diagram{\begin{tikzcd}[ampersand replacement=\&,column sep=large]
\& \text{mult. fini} \\
(\mu_p) \arrow[ur,dash] \arrow[dr,dash] \& \\
\& \text{connexe}
\end{tikzcd}}

\diagram{\begin{tikzcd}[ampersand replacement=\&,column sep=large,row sep=large]
{} \arrow[r,dash] \& {} \arrow[dr,dash] \& \text{\uline{\unsure{Unip.t conn.}}} \arrow[dr,dash,"\mathbb Z/\ell\mathbb Z"] \& \\
\& \uline{G_m} \arrow[ur,dash,"\mu_p"] \arrow[dr,dash,"\mathbb Z/\ell\mathbb Z"']
\& {} \arrow[dr,dash,"\mathbb Z/\ell\mathbb Z"] \& \uline{\mathbb Z/p\mathbb Z} \\
\& \& {} \arrow[ur,dash,"\mu_p"'] \& \text{\uline{\unsure{Unip.t connexe}}} \arrow[uu,dash]
\end{tikzcd}}
\scanpage{181}
% archive sheet number: 180

Modules discrets \deleted{\unsure{sur}} de type fini sur
\[
A=\deleted{\unsure{\mathbb W|W}\times}\lim_{\longleftarrow\,n}\mathbb W_{n\,\sigma}[F,\unsure{V}]/(V^n,FV-p)
\]

(Catégorie $C_1$) on considère $C_1^\circ=C_2$.

Et la sous-catégorie \inserted{pleine} $C'_1$ de $C_1^{\unsure{\alpha|\circ}}$ formée des modules locaux annulés par une puissance de $F$, [qui en $C_1^\circ$ correspond \unsure{aux} \deleted{\unsure{limites projectives}} \inserted{à la sous-catégorie $C'_2$} des groupes finis \uline{radiciels} unipotents].

\margin{$C_1$ \deleted{Modules de}\par $C_2$}
\margin{$C'_1$\par $\unsure{\wr|21}$\par $C'_2$}

La dualité de Cartier des groupes finis, définie \deleted{par} sur les modules de type \unsure{cocycliques|cocyclique} définit un \uline{isomorphisme} $C'_1\simeq C'_2$.
D'ailleurs, \unsure{tt objet} $G$ de $C_2$ a un sous-objet canonique tel que $G/G'$ soit le plus grand quotient dans $C'_2$, et de $=$ tout objet de $C_1$ a un plus grand sous-objet \deleted{\unsure{en plus}} dans $C'_1$. Un objet de la nouvelle catégorie $C$ serait un \deleted{\unsure{couple}} \inserted{triple} formé de \deleted{\unsure{un}} $(\deleted{\unsure{Q}}X,Y,u)$, où $X\in C_2$, $Y\in C_1$, et $u$ est un isom de $\alpha(X)$ sur $\beta(Y)$.

\diagram{\begin{tikzcd}[ampersand replacement=\&]
C_1 \arrow[dr,"\alpha"] \& \\
\& C_0 \\
C_2 \arrow[ur,"\beta"'] \&
\end{tikzcd}}
\scanpage{182}
% archive sheet number: 181

\[\underset{\unsure{m|\mathfrak m}}{A}\qquad\deleted{\unsure{M}}\]

$X,Y$ deux modules sur $A$
\[H^0_m(X)\xrightarrow{\sim}DH^0_m(Y)\]

$(X,Y)\unsure{\longleftarrow|\longrightarrow}(X',Y')$ consiste en
\[\begin{aligned}X&\longrightarrow X'\\Y&\longleftarrow Y'\end{aligned}\]

tel que le diagramme
\diagram{\begin{tikzcd}[ampersand replacement=\&]
H^0(X) \arrow[r] \arrow[d,"\sim"'] \& H^0(X') \arrow[d,"\sim"] \\
DH^0(Y) \arrow[r,dash] \& DH^0(Y')
\end{tikzcd}}
soit \uuline{commutatif}

Si on veut inclure les groupes $\unsure{G_m|\mathbb G_m}$ (\deleted{\unsure{dans les}} parmi les algébriques (avec aussi $\mu_p,\mathbb Z/\ell\mathbb Z$) et $\mathbb Z$ parmi les groupes formels (\unsure{bien aussi|avec aussi} $\mathbb Z/p\mathbb Z,\mathbb Z/\ell\mathbb Z$), alors la \unsure{prés.} est la $=$, en prenant comme nouvel anneau topologique
\[\mathbb Z\times A\]
et le $H^0$ étant relatif aux idéaux $m\unsure{=|:}(F,V)$ de $A$ et les idéaux non nuls quelconques de $\mathbb Z$, la dualité
\scanpage{183}
% archive sheet number: 182

\unsure{mélangera|mélangeant} ces deux \uuline{facteurs}

Dans une catégorie de groupes comme dessus, le quotient est \unsure{tjs|tjrs} annulé par $F^{\unsure{\infty|n}}$, et le sous groupe par $V^\infty$. Mais \unsure{passons|prenons} aux ind-objets de $C_1$ et aux pro-objets de $C_2$, les uns ni les autres ne sont annulés par $F^n$ resp. $V^n$. Comment définir le \unsure{recollement} maintenant ??

Quelle est l'interprétation des modules de type fini sur
\[\unsure{\mathbb W|W}_{\infty\,\sigma}[F,V]/(FV\unsure{-}p)\qquad ??\]

\unsure{Je pense} que c'est l'anneau des endomorphismes de $\mathbb W_\infty$ qui \unsure{invarient} le \unsure{ss. grpe|sous groupe} $D(\mathbb W_\infty)$ [\deleted{\unsure{f|les}} vecteurs de Witt dont toutes les composantes sont des nilpotents, et tous sauf un nb fini sont nuls ---], ou les endomorphismes de $D(\mathbb W_\infty)$ qui \unsure{invarient} l'accouplement fondamental de ce \unsure{grpe|groupe} \unsure{avec|sur} lui-$=$.
\scanpage{184}
% archive sheet number: 183
% Large cancellation strokes cross the left notes and the central and lower-left blocks.

\[\deleted{\unsure{G_m|\mathbb G_m}\;\text{---}\;\unsure{\operatorname{Ext}}}\]
\[\unsure{\mathbb W|W}_\sigma[F,V]/(V^n,FV-p)\]

\diagram{\begin{tikzcd}[ampersand replacement=\&]
\deleted{0} \arrow[r,dash] \& \deleted{\unsure{\alpha}\mu_n} \arrow[r,dash] \& \deleted{G_m} \arrow[r,"{\deleted{n}}"] \& \deleted{G_m} \arrow[r,dash] \& \deleted{0}
\end{tikzcd}}
\[\deleted{\unsure{V}}\quad\deleted{V^n}\quad k[V]\]
\diagram{\begin{tikzcd}[ampersand replacement=\&,column sep=small]
0 \arrow[r,dash] \& \mathbb Z \arrow[r,dash,"n"] \& \mathbb Z \arrow[r,dash] \& \operatorname{Ext}^{\unsure{1}}(G_m,\mu_n) \arrow[r] \& 0
\end{tikzcd}}
\[k_\sigma[F]\]

\[\deleted{\widetilde{\mathbb W}_\infty}\]
\deleted{\unsure{mod.}}\qquad \deleted{artinien}

\deleted{\uline{pr. artinien} et \uline{noethérien}}
\[k[F]\qquad\deleted{\unsure{F\mathbb Z_F}}\quad k[F]\]

\diagram{\fbox{\begin{tikzcd}[ampersand replacement=\&,column sep=small,row sep=large]
\& \deleted{\mu_p} \arrow[rrd,dash,"{\deleted{\alpha_p}}"] \& \& \\
\deleted{\uline{G_m\,;\,G_a}} \arrow[ur,dash,"{\deleted{\mathbb Z/\ell\mathbb Z}}"]
\arrow[rrrd,dash,"{\deleted{\mu_p\quad\alpha_p\quad\mathbb Z/p\mathbb Z}}"']
\& \& \& \deleted{\mathbb Z/p\mathbb Z} \& \deleted{\uline{\unsure{\mathbb V|V}}}\quad\deleted{\uline{\mathbb Z}} \\
\& \& \& {} \arrow[u,dash,"{\deleted{\mathbb Z/\ell\mathbb Z}}"'] \&
\end{tikzcd}}}

\diagram{\begin{tikzcd}[ampersand replacement=\&]
\deleted{\mathbb W_\sigma\{F,V\}/(FV-p)} \arrow[d] \\
\deleted{\mathbb W_\sigma[[F,V]]/(FV-p)}
\end{tikzcd}}

(de dimension 2)

\uuline{Anneau de Cartier}

\margin{\unsure{Classifie les pro-groupes algébriques unipotents} \deleted{\unsure{conn.}} ou les groupes \deleted{\unsure{formels}} (\uline{ind finis})}

\deleted{\unsure{formels infinitésimaux}}

\uuline{Anneau de Gabriel-Dieudonné}

complété du précédent pour l'idéal \inserted{maximal} $(F,V)$

\deleted{Classification des}
\scanpage{185}
% archive sheet number: 184

\[\begin{aligned}
\unsure{\mathbb W|W}_\infty&\;\text{---}\; \mathbb W_\infty\\
&=\lim_{\longrightarrow}\operatorname{Hom}(\mathbb W_\infty,\mathbb W_{\unsure{n|{-}}})
\end{aligned}\]
\[\deleted{\operatorname{Hom}}D(\mathbb W_\infty)\subset\overline{\mathbb W}_\infty\quad\unsure{1}\]
\[\deleted{\unsure{\mathbb Z|D}}\]

\[D(\mathbb W_\infty)\qquad F\]
\[D(W_\infty)\longrightarrow W_\infty\]
\[D(\uline{W_\infty})\times D(\uline{W_\infty})\]
\[\operatorname{Hom}(\mathbb W_\infty,\unsure{G_m|\mathbb G_m})\]
\[\mathbb W_\infty\times\]

\[
\left\{\begin{aligned}
&\text{groupes algébriques affines \inserted{comm}}\ =\
\text{algèbres de type fini sur $k$,}\\
&\hspace{3em}\text{avec $\Delta$ \unsure{co-comm.}}
\left\{\begin{array}{l}\text{\unsure{unitaire}}\\\text{\unsure{inversible}}\end{array}\right.\\[1em]
&\text{\deleted{\unsure{gr}} groupes ($\deleted{\unsure{gr}}\mathrm{pro}$ (algébriques affines))}\ =\
\text{\deleted{\unsure{Alg.}} \unsure{les $=$ sans, les type fini|les $=$, sans le type fini}}\\[1em]
&\left\{\begin{aligned}
&\text{\uline{Question} est-ce un \deleted{\unsure{un}} pro (gr alg. affine)}\\
&\text{i.e. l'algèbre est-elle limite inductive d'algèbres de type fini}\\
&\text{qui sont stables par $\Delta$ ??}
\end{aligned}\right.
\end{aligned}\right.
\]
\scanpage{186}
% archive sheet number: 185

\deleted{\unsure{coalgèbres|cogèbres}} \deleted{groupes} objets
\[
\text{(ind (algébriques finis))}\simeq\text{coalgèbres comm \unsure{nn. counitaires}}
\]
\margin{car \deleted{\unsure{on}} on est sur un corps, et on peut utiliser la dualité}
\deleted{\unsure{les}} \uline{engendrées par} des \uline{coalgèbres de} \uline{rang fini sur $k$}

\deleted{[une telle coalgèbre est en effet engendrée par des coalgèbres}
(cette dernière condition sans doute superflue, i.e. une multiplication \inserted{\unsure{comm non unitaire}} \uline{continue} sur un $k^I$ \deleted{\unsure{n'est}} \inserted{\unsure{n'impose|n'implique}} pas nécessairement une \deleted{\unsure{système}} \unsure{limite} projective d'algèbres de rang fini sur $k$)

\margin{$\Delta(u)\in U\otimes U)$}
\diagram{\begin{tikzcd}[ampersand replacement=\&]
U \arrow[dr,leftrightarrow] \& \unsure{V|U} \arrow[dl,leftrightarrow] \\
V^* \& U^*
\end{tikzcd}}

\deleted{\unsure{gr}}
\[
\underbrace{\text{groupe (ind (algébriques finis))}}_{\text{\tikz[baseline]{\draw (0,0) ellipse (1.35cm and .30cm) node {groupe formel};}}}
\simeq\text{coalgèbre du type précédent,}
\]
munie d'une loi d'algèbre ass. commutative unitaire compatible avec la précédente

\margin{\uline{Exemple} l'Algèbre affine d'un groupe algébrique linéaire ou celle d'un pro. (groupe algébrique linéaire)}

d'après Cartier, \inserted{en car. $p$} une telle algèbre serait limite inductive d'algèbres de \inserted{\unsure{ce type, qui sont}} de type fini du pt de vue associatif [du moins sur un corps alg. clos, où la partie discrète se sépare de la partie infinitésimale]

Donc il en résulterait une dualité entre
\[
\mathrm{pro}\ (\text{grps algébriques affines})\sim\mathrm{grp}\ (\text{ind. Alg. finis}).
\]
Le plus beau serait \unsure{maintenant|naturellement} qu'on ait le résultat que \unsure{cela} correspond aux \unsure{homomorphismes|isomorphismes} \uline{\unsure{quelconques|algébriques}}
\scanpage{187}
% archive sheet number: 186

La \unsure{catégorie|donnée} d'un \deleted{\unsure{tel}} groupe formel ou groupe pro-algébrique affine \inserted{\unsure{définis comme certains groupes} (ind ($k$-affines))} \unsure{définissent|définit} donc une \deleted{coalgèbre} hyperalgèbre topologique sur $k$, dont l'espace sous-jacent est linéairement compact. Il n'est pas exclu que \unsure{les|ttes} hyperalgèbres topologiques ayant cette propriété \deleted{\unsure{puissent}} puissent être interprétées ainsi. \deleted{\illegible{2}}
% slanted insertion at the left, linked to the word affine
% slanted
\margin{\deleted{\unsure{on peut vérifier|qui par} \unsure{que|des filtres} \illegible{4} \unsure{la} \illegible{2}}}

Le mécanisme de dualité de Cartier est alors immédiat, ainsi que son interprétation en termes d'accouplement dans $G_m$. La définition de $F$ et $V$ est aussi immédiate, en car. $p$.

% Sweeping cancellation strokes cross the lower comparison.
\diagram{\begin{tikzcd}[ampersand replacement=\&,column sep=large,row sep=small]
\deleted{\unsure{\mathbb G_m|G_m}} \& \deleted{\mathbb Z/\ell} \& \mu_p \\
\deleted{\unsure{F=x^p}} \& \& \\
\deleted{F(x)=x^p} \& \deleted{F(x)=x^{\deleted{\unsure{p}}}} \& \deleted{F=0} \\
\deleted{V(x)=x} \& \deleted{V(x)=x^p} \& V=\mathrm{id} \\
\& \text{($F$ et $V$ \unsure{absolus|disons})} \& \\
% slanted note, linked to both F formulas
\& \text{\inserted{\unsure{par la réduction naturelle au schéma} : $\mathbb F_p$ \unsure{elle identifie} $\mathbb G_m=\mathbb G_m^{(p)}$}} \arrow[uuul] \arrow[uuu] \&
\end{tikzcd}}
\scanpage{188}
% archive sheet number: 187

$A$ algèbre affine d'un groupe affine $G$ (pas néc. de type fini) sur l'anneau $k$

\[G\unsure{\longrightarrow|\text{---}}\unsure{G_m|\mathbb G_m}\quad\text{hom de $k$-grp}\quad\unsure{\approx|\simeq}\quad G\longrightarrow\mathbb G_a\]
morphisme de $k$-schémas, tel que
\begin{enumerate}[label=(\roman*)]
\item Addition transformée en multiplication
\item $0$ s'en va en $1$
\end{enumerate}

i.e. $f\in A$, telle que
\[\Delta f=f\otimes f\tag*{a)}\]
\[\varepsilon(f\unsure{)|\}}=1\tag*{b)}\]

Supposant que $A$ soit \uline{libre}, alors $A$ est l'ensemble des formes $k$-linéaires continues sur $\check A=\operatorname{Hom}(A,k)$, muni de la convergence simple. Alors a) signifie que $f$ est multiplicatif \deleted{\unsure{et}} [que $\check A$ est muni de la structure de $k$-algèbre déduite de $\Delta$] et b) que $f$ est unitaire; donc \unsure{a/. et b|c/. et b}, que $f$ est un hom d'algèbres topologiques.

\deleted{D'où, si} Note que rien ne garantit que $\check A$ soit topologisé linéairement en tant qu'anneau, i.e. qu'il y ait un système fondamental de voisinages qui soient des \uline{idéaux}; \unsure{or|ou}, ce qui revient au $=$ si $k$ est un corps [et ce qui suffit en tous cas]. rien ne
\scanpage{189}
% archive sheet number: 188

\unsure{garantit} que $A$, en tant que co-algèbre, soit limite inductive de ``ss-Coalgèbres'', de type fini en tant que modules sur $k$ \deleted{\unsure{s'il en est ainsi}} \deleted{et, de préférence} [et de préférence facteurs directs en tant que modules sur $k$]. \uline{S'il} en est cependant ainsi, alors $\check A$ peut être regardé comme l'algèbre affine d'un ind-(schéma affine) sur $k$, \inserted{[pas nécessairement un schéma formel]} \deleted{\unsure{, qui sera un groupe}} et \deleted{sera}\inserted{définit} \unsure{=} un \uline{groupe} (ind-affine) sur $k$, désigné alors par $D(G)$ et appelé \uline{groupe dual} (au sens Cartier) de $G$. \deleted{\unsure{Ce sera le}}\inserted{Le foncteur} \inserted{de $k$-algèbres} représenté par ce \deleted{\unsure{progr.}} groupe ind.affine sera précisément $\underline{\mathrm{Hom}}_{k\text{-grps}}(G,\mathbb G_m)$.
% slanted
\margin{Si, \unsure{dans le} cas d'un corps cela \unsure{va de bien|va de soi}\unsure{...}}

Si \deleted{$G$} $k$ est \deleted{de type fini} un corps et $G$ de type fini sur $k$, alors si $k$ est de car $p>0$, d'après Cartier $D(G)$ est bien un groupe (ind-fini).
\[
\left\{\begin{aligned}
D(\mathbb G_m)&=\deleted{\unsure{\mathbb Z}}\mathbb Z\\
D(\mathbb Z/\ell\mathbb Z)&\simeq\mathbb Z/\ell\mathbb Z
\quad\text{\deleted{\unsure{naturellement}} si $\ell$ premier \unsure{à car.}, et}\\[-2pt]
&\hspace{28mm}\text{si les racines primitives $\ell$-ièmes de l'unité \unsure{existent}}\\[-2pt]
&\hspace{28mm}\text{dans \deleted{\unsure{$\Omega$}}$k$}\\
D(\mathbb Z/p^n\mathbb Z)&\simeq\mu_{p^n}\quad\text{\deleted{\unsure{par ce}}}\\
D(\mu_{p^n})&\unsure{=|\simeq}\mathbb Z/p^n\mathbb Z\\
D(\mathbb G_a)&=\Bigl(\text{Algèbre enveloppante \inserted{\uline{\unsure{avec distinction de puiss $p^i$|sans distinction de puiss $p^i$|avec définition de puiss $p^i$}}}}\\[-2pt]
&\hspace{8mm}\text{de la $p$-algèbre de Lie libre engendrée par un élément}\\[-2pt]
&\hspace{8mm}\text{(sur un corps de car $p>0$)}
\end{aligned}\right.
\]
\scanpage{190}
% archive sheet number: 189

Soit $\underline C$ une catégorie, \unsure{déf:|d'où}
\[
\operatorname{Pro}(\underline C)\unsure{\subset|=}\underline{\mathrm{Hom}}(\underline C,(\mathrm{Ens}))
\]
Un objet \inserted{$\unsure{\mathscr X|\mathcal X}$} de $\operatorname{Pro}(\underline C)$ est connu (à \unsure{isomorphisme} près) quand on connait le foncteur $\unsure{\mathscr Y|\mathcal Y}\mapsto\mathrm{Hom}_{\operatorname{Pro}(\underline C)}(\mathscr Y,\mathscr X)$ de $\operatorname{Pro}(\underline C)^\circ$ dans $(\mathrm{Ens})$, \deleted{\unsure{mais}} mais,

\deleted{$\mathscr Y=(Y_j)_{j\in\unsure{I|J}}$, les $Y_i\in\underline C$, alors \unsure{...},}
\[
\deleted{\mathrm{Hom}_{\operatorname{Pro}(\underline C)}(\mathscr Y,\mathscr X)=\varprojlim_{\unsure{i|j}}\mathrm{Hom}(\mathscr Y,\mathscr X_i)=}
\]
identifiant $\underline C$ à une sous-catégorie de $\operatorname{Pro}(\underline C)$, on \unsure{notera|voit} que $\mathscr X$ n'est \uline{plus connu} à isom près quand on connait le foncteur $Y\mapsto\mathrm{Hom}_{\operatorname{Pro}(\underline C)}(Y,\mathscr X)$ de $\underline C^\circ$ dans $(\mathrm{Ens})$. [Ex: $\underline C=(\mathrm{Ens})$, alors $(X_i)$ et $X=\varprojlim_{i\,(\mathrm{Ens})}X_i$ définissent le $=$ foncteur $\underline C^\circ\mathbin{-}(\mathrm{Ens})$, mais en général ne sont pas isomorphes ds $\operatorname{Pro}(\underline C)$]. Le foncteur naturel
\[
\operatorname{Pro}(\underline C)\longrightarrow\underline{\mathrm{Hom}}(\underline C^\circ,(\mathrm{Ens}))
\]
n'est \inserted{donc} \uline{pas} pleinement fidèle en général. Il importe de bien distinguer en particulier, si $X=\varprojlim_{\underline C}\mathscr X$ existe, entre $\mathscr X$ et $X$

[on peut dire seulement qu'on a un morphisme naturel ds $\operatorname{Pro}\underline C$
\[
X=\varprojlim_{\underline C}\mathscr X\longrightarrow\mathscr X
\]
\scanpage{191}
% archive sheet number: 190

universel pour les \inserted{$\operatorname{Pro}(\underline C)$-}morphismes d'objets de $\underline C$ en $\unsure{\mathscr X|\mathcal X}$; si $\varprojlim_{\underline C}\mathscr X$ existe pour tt $\mathscr X$, cela signifie précisément que le foncteur \deleted{\unsure{d'inclusion}} pleinement fidèle $\underline C\longrightarrow\operatorname{Pro}(\underline C)$ admet un \deleted{\unsure{droit}}\inserted{adjoint} de Kan, ``projection'' de $\operatorname{Pro}(\underline C)$ \deleted{\unsure{sur}} sur $\underline C$.

On ne confondra pas cette situation avec la suivante: \deleted{\unsure{alors}} il n'y a pas néc. de limites projectives ds $\underline C$, mais on a plongé $\underline C$ comme \uline{sous-catégorie pleine} ds $\underline C'$ telle que \deleted{\unsure{tt|tout}} système projectif de $\underline C$ ait une limite projective \uline{ds $\underline C'$}; alors on a un foncteur naturel
\[
\operatorname{Pro}(\underline C)\longrightarrow\underline C'
\]
qui \uline{parfois} \deleted{\unsure{défini}} (en un cas fréquent) sera pleinement fidèle. [Il en est p.ex. ainsi \unsure{chaque} fois que les objets de $\underline C$ sont artiniens ??]

\uline{Exemple} les pro-(schémas affines de type fini) sur $S$ \deleted{\unsure{soit}} (schéma quasi-compact ou \unsure{préschéma} loc.\unsure{noeth\textsuperscript{n}}) s'identifient exactement aux \uline{schémas} \uline{affines} sur $S$.
% slanted
\margin{N.B. $\operatorname{Pro}(\underline C\text{-grps})\longrightarrow(\operatorname{Pro}\underline C)\text{-grps}$ est un foncteur \uline{pleinement} \uline{fidèle}}

\uline{Question} \deleted{Supposons $S$ \unsure{noethérien}.} Un schéma en groupes affine sur $S$ est-il un pro (groupe affine de type fini), i.e. si $S=\operatorname{Spec}(k)$
\scanpage{192}
% archive sheet number: 191

($k$ un anneau) est-il vrai que si $A$ est une hyperalgèbre (\unsure{commutative?|non commutative?}) sur $k$, alors $A$ \deleted{\unsure{commut.}} est limite inductive de \deleted{\unsure{sous-algèbres de}}\deleted{\unsure{hyperalgèbres}}\inserted{hyperalgèbres} \uline{de type fini} \deleted{\illegible{3}}\inserted{\unsure{du pt de vue associatif}}. Il suffit pour ceci que $A$ soit limite inductive de ss-$k$-modules de type fini $V$ tels que $\Delta V\subset\operatorname{Im}V\otimes V$, et de plus [$A_V$ dés la ss-$k$-algèbre associative engendrée] \deleted{\unsure{et}} $A_V\otimes A_V\mathbin{-}A\otimes A$ soit \uline{injective}. \deleted{\unsure{Ce qui}}
\margin{\diagram{\begin{tikzcd}[ampersand replacement=\&,row sep=small]
G\arrow[r,"\text{épim}"]\&G'\\
A\&A'\arrow[l]
\end{tikzcd}}}

On voit facilement qu'il en est ainsi si $k$ est un \uline{corps} [alors la \deleted{première}\inserted{deuxième} \unsure{cond\textsuperscript{on}} est remplie automatiquement; pour la première, il est possible qu'elle soit vraie \deleted{\unsure{au moins}}\inserted{sous des} \deleted{si $k$ est noethérien:} conditions beaucoup plus générales, telles que: $A$ libre].
\margin{$\unsure{\ast|+}$}

\uline{Sur un corps} $k$, \deleted{\unsure{un groupe affine}}
\[
\underline{\text{pro (groupe alg. affine)}}=\underline{\text{grp. \unsure{alg.|sch.} affine}}
\]
\deleted{\unsure{$G$}}

Le ``dual'' d'un \deleted{\unsure{tel}} groupe \inserted{comm. sur $k$} satisfaisant la condition (1) est un groupe (ind.fini) $\deleted{\unsure{D}}\unsure{\widetilde G|\widehat G}$, sur $S$, dont l'hyperalgèbre est l'algèbre affine de $G$. \unsure{Se bornant au} type fini sur un corps alg. clos, on trouve encore
% slanted
\margin{\unsure{groupe formel} au sens de Cartier}
\scanpage{193}
% archive sheet number: 192

types \deleted{simples}\inserted{``élémentaires''}, de part et d'autre,
\[
\unsure{\mathbb G_m|G_m},\quad\mathbb Z/\ell\mathbb Z,\quad\mu_p\quad\unsure{\mathbb G_a|G_a},\quad\alpha_p,\quad\mathbb Z/p\mathbb Z
\]
\margin{$\mu_p,\alpha_p,\mathbb Z/p\mathbb Z$ \unsure{spécial à} car. $p>0$}
[Et, l'ordre où on les rencontre dans une suite de composition canonique \unsure{est} $=$, à très peu de chose près, sur un corps non alg. clos ni $=$ parfait]
% slanted
\margin{Si on se borne à l'unipotent fini, on le \unsure{trouve en prenant|trouve en faisant} \unsure{relativement|des extensions} en bas [le \uline{multiplicatif} \unsure{correspond aux} discrets séparables]. Or Cartier a \deleted{\unsure{exactement}} classifié les \deleted{hyperalgèbres}\inserted{\unsure{grp. alg. finis}} \uline{unipotents}}
\[
\mu_p,\quad\alpha_p,\quad\mathbb V,\quad\mathbb Z/p\mathbb Z,\quad\mathbb Z/\ell\mathbb Z,\quad\mathbb Z
\]
$\mathbb V=D(\mathbb G_a)$ est le \unsure{cinquième} groupe; en car $p>0$, on a envie de les mettre ensemble en faisant des \uline{extensions}, \deleted{\unsure{Prenons une}}\inserted{dans la} \deleted{\unsure{propre}} catégorie (présumée abélienne) des groupes (ind-affines) \deleted{[\unsure{formels}]} [ind-$S$-schémas affines, avec des structures dessus]
% slanted
\margin{\unsure{Regarder} en général (sans condition de commutativité), \deleted{\unsure{pro-catégorie}} des ``groupes \deleted{algébriques} proalgébriques'' (extension d'un groupe formel par un \unsure{pro-grpe algébrique})]}

Un tel groupe aura une algèbre affine, qui sera une \deleted{limite projective \unsure{stricte}}\inserted{somme directe \inserted{infinie}} de \unsure{produits} de \unsure{discrets} (\unsure{Bo---|Bour---}) \deleted{\unsure{par}} (c'est un \unsure{$\beta$|$f$} \unsure{venant} d'un dans une \unsure{limite}, \unsure{un peu} dans le cas non \unsure{proalgébrique})

\deleted{\unsure{Opération} de $F,V$},
% slanted
\margin{dans le cas $k$ parfait, extensions triviales, \unsure{du moins} \unsure{le sont pour} \illegible{2}}
\diagram{\[\begin{array}{c@{\qquad}c@{\qquad}c}
\text{Multiplicatif [\unsure{quelle compt. est. connexe})}
&\text{unipot\deleted{\unsure{ent}}}
&\text{[\unsure{sans $\ell$-torsion ou $p$-torsion}]}\\
&\vert\!\rule{34mm}{0.5pt}\!\vert
&\text{groupes discrets \unsure{ét.}}\\
\text{Tore}\;\backslash\;\text{grpe profini \deleted{\unsure{étale}}}\;\times\mu_p
&&\vert\!\rule{40mm}{0.4pt}\!\vert\\
\text{\hspace{15mm}d'ordre premier à $p$}
&\text{\deleted{\unsure{groupe}} admet un}
&\text{grpe \deleted{\unsure{de torsion}}\inserted{\unsure{ordinaire}}}\\
\text{\unsure{dual aux} groupes \deleted{\unsure{libres}}\inserted{\unsure{discrets}}}
&\text{plus grand sous grpe alg. et}
&\text{\deleted{\unsure{à $\pi$ opère}}}\\
\text{\unsure{sans torsion à}\;$\unsure{\pi}$}
&\text{un plus gd \deleted{\unsure{partie}}\inserted{\unsure{quotient}}}
&\text{\unsure{où} $\pi$ opère}\\
\text{$\vert$ \unsure{dual aux} groupes de torsion première à $p$}
&\text{infinitésimal (sur $k$)}
&\text{continûment}\\
\text{\unsure{où opère} $\pi$ $\vert$ \unsure{dual aux grpes}}&&\\
\text{\unsure{de $p$-torsion où opère} $\pi$}&&
\end{array}\]}
